\documentclass[12pt,a4paper]{article}
\usepackage[utf8]{inputenc}
\usepackage[T2A]{fontenc}
\usepackage[russian]{babel}
\usepackage{amsmath,amssymb}
\usepackage{booktabs,array}
\usepackage{geometry}
\usepackage{microtype}
\usepackage{xurl}
\geometry{margin=2.3cm}

\title{Preservation-аудит проекта\\
после Wilson modular-state gate}
\author{}
\date{6 августа 2026 г.}

\begin{document}
\maketitle

\section{Проверенная задача}

Аудит отвечает не на вопрос о физической истинности формул, а на более
технический вопрос: не были ли прежние результаты потеряны, перезаписаны или
сломаны при добавлении новых отрицательных gate-тестов.

\section{Результат}

Основные документы сохранены как крупные самостоятельные корпуса:
\begin{center}
\begin{tabular}{lrr}
\toprule
Файл & Размер & Собранный PDF \\
\midrule
\texttt{main.tex} & $510839$ байт & $112$ страниц \\
\texttt{tome2\_s2t\_spectral\_closure.tex} & $539046$ байт & $165$ страниц \\
\bottomrule
\end{tabular}
\end{center}

Во втором томе по-прежнему присутствуют блоки $S_{vac}$, $m_\tau$,
$23+\pi^{-1}$, C6, SU(5), parent action и глобальный falsification-аудит.
В тот же основной файл добавлена часть II.B с family-selector,
inverse-susceptibility и новым Definition of Done.
В первом томе сохранены послесловие после проверки вторым томом, разделение
математического ядра, условных связей и отрицательно закрытых ветвей, а также
перспектива следующего тома.

Структурные проверки дали:
\begin{itemize}
 \item по одному \texttt{\string\end\{document\}} в каждом главном файле;
 \item баланс ключевых LaTeX-окружений;
 \item отсутствие удалённых Git-путей;
 \item наличие новых family-selector, CKM, \(C_4\), chiral, zero-prompt и observed-world scripts/results/gates;
 \item успешную синтаксическую компиляцию $152$ audit-скриптов;
 \item успешный разбор $167$ JSON-файлов результатов;
 \item полную сборку обоих главных томов.
\end{itemize}

\section{Статус формул}

Новые gate-тесты не удаляли и не переписывали прежние формулы. Они добавлялись
отдельными документами и меняли только эпистемический статус некоторых
механизмов:
\begin{equation}
 \text{сохранённая формула}
 \quad\neq\quad
 \text{доказанный операторный вывод}.
\end{equation}
Численные соотношения, геометрические конструкции, нейтринный submodel,
спектральные результаты и все ранее сформулированные no-go остаются в корпусе.
Закрыты лишь конкретные попытки вывести их из недостаточно богатого parent
action или intrinsic family-selector текущего носителя.

Машинный отчёт сохранён в
\path{s2t_project_preservation_audit.py} и
\path{s2t_project_preservation_results.json}.

\end{document}